
\documentclass[11pt,a4paper]{article}
\usepackage[utf8]{inputenc}
\usepackage[T1]{fontenc}
\usepackage[spanish,es-noshorthands]{babel}
\usepackage[margin=2.2cm]{geometry}
\usepackage{booktabs,amsmath,amssymb,microtype,xcolor}
\usepackage[colorlinks=true,linkcolor=black,urlcolor=blue]{hyperref}
\usepackage{titlesec}
\titleformat{\section}{\large\bfseries}{\thesection.}{0.6em}{}
\titleformat{\subsection}{\normalsize\bfseries}{\thesubsection.}{0.6em}{}
\setlength{\parskip}{4pt}
\setlength{\parindent}{0pt}
\renewcommand{\arraystretch}{1.12}

\title{\vspace{-1.2cm}\textbf{Qué ve el kernel}\\[2pt]
\large La representación de la fuente con y sin BatchNorm,\\
\large y por qué decide el resultado de la adaptación}
\author{}
\date{}

\begin{document}
\maketitle
\vspace{-1.4cm}

\begin{center}
\fbox{\begin{minipage}{0.93\textwidth}\small
\textbf{En una línea.} Sin BatchNorm el codificador de la fuente entrega al kernel
un espacio de \emph{una} dirección efectiva para un problema de \emph{tres} clases.
Tres clases sobre una recta no se pueden separar con ningún kernel de distancia:
la del medio queda entre las otras dos por construcción. No es un problema de
sintonía --- es geometría, y se resuelve antes de que al kernel se le pregunte nada.
\end{minipage}}
\end{center}

\section{Las tablas, primero}

Todas las medidas salen de \textbf{la fuente sola}, al terminar la fase 1: mismo
seed, mismos datos, mismo orden de lote. Lo único distinto entre las dos columnas
es BatchNorm. Las dos alcanzan \textbf{1.0000} de acierto sobre la fuente, así
que lo que difiere no es el ajuste sino la forma de lo que se le entrega al kernel.

\subsection*{Sin el término KL}
\noindent
\begin{minipage}[t]{0.49\textwidth}
  \centering
  \footnotesize\setlength{\tabcolsep}{3pt}
  \begin{tabular}{lrrrr}
      \multicolumn{5}{c}{\textbf{SIN BatchNorm}} \\[2pt]
      \toprule
      familia & dim.\,ef. & \% entre & contraste & acierto \\
      \midrule
      rollo & 1.97 & 92.3\% & 0.4383 & 0.9401 $\pm$ 0.109 \\
      toroides & 1.11 & 95.5\% & 0.5536 & 0.3651 $\pm$ 0.030 \\
      lunas & 1.04 & 99.9\% & 0.5799 & 0.3306 $\pm$ 0.006 \\
      \bottomrule
    \end{tabular}
\end{minipage}\hfill
\begin{minipage}[t]{0.49\textwidth}
  \centering
  \footnotesize\setlength{\tabcolsep}{3pt}
  \begin{tabular}{lrrrr}
      \multicolumn{5}{c}{\textbf{CON BatchNorm}} \\[2pt]
      \toprule
      familia & dim.\,ef. & \% entre & contraste & acierto \\
      \midrule
      rollo & 1.97 & 99.1\% & 0.4246 & 0.6769 $\pm$ 0.014 \\
      toroides & 1.86 & 91.3\% & 0.4757 & 0.6299 $\pm$ 0.098 \\
      lunas & 1.95 & 99.8\% & 0.4228 & 0.4988 $\pm$ 0.133 \\
      \bottomrule
    \end{tabular}
\end{minipage}
\vspace{11pt}


\subsection*{Con el término KL}
\noindent
\begin{minipage}[t]{0.49\textwidth}
  \centering
  \footnotesize\setlength{\tabcolsep}{3pt}
  \begin{tabular}{lrrrr}
      \multicolumn{5}{c}{\textbf{SIN BatchNorm, con KL}} \\[2pt]
      \toprule
      familia & dim.\,ef. & \% entre & contraste & acierto \\
      \midrule
      rollo & 2.81 & 81.9\% & 0.3580 & 0.7355 $\pm$ 0.251 \\
      toroides & 2.83 & 80.7\% & 0.3622 & 0.2781 $\pm$ 0.068 \\
      lunas & 2.33 & 90.5\% & 0.3946 & 0.3346 $\pm$ 0.006 \\
      \bottomrule
    \end{tabular}
\end{minipage}\hfill
\begin{minipage}[t]{0.49\textwidth}
  \centering
  \footnotesize\setlength{\tabcolsep}{3pt}
  \begin{tabular}{lrrrr}
      \multicolumn{5}{c}{\textbf{CON BatchNorm, con KL}} \\[2pt]
      \toprule
      familia & dim.\,ef. & \% entre & contraste & acierto \\
      \midrule
      rollo & 2.80 & 82.0\% & 0.3595 & 0.9451 $\pm$ 0.079 \\
      toroides & 2.83 & 81.3\% & 0.3606 & 0.7657 $\pm$ 0.206 \\
      lunas & 2.83 & 80.8\% & 0.3605 & 0.7062 $\pm$ 0.109 \\
      \bottomrule
    \end{tabular}
\end{minipage}
\vspace{11pt}


\textit{dim.\,ef.} son las direcciones efectivas (de 1 a 16); \textit{\% entre} es
la parte de la varianza que es separación entre clases; \textit{contraste} es
$\bar K_{=} - \bar K_{\neq}$; \textit{acierto} es sobre el dominio objetivo, 5
semillas. El azar es $0.3333$.

\section{Qué experimento se hizo}

Tres familias de juguete (\textit{rollo}, \textit{toroides}, \textit{lunas}),
reducidas a \textbf{3 clases}. Cada una tiene un dominio fuente y uno objetivo; el
objetivo es el mismo generador con una traslación de $5{,}0$ por coordenada, una
rotación de $45^\circ$ y ruido $\sigma = 0{,}3$.

\textbf{El montaje son dos fases.} La fase 1 entrena \emph{un} codificador
$f_\theta: \mathbb{R}^d \to \mathbb{R}^{16}$ y una cabeza lineal sobre la fuente,
con entropía cruzada, y \textbf{congela los dos}. Una copia única siembra el
codificador objetivo $f_{\theta'}$, que de ahí en más entrena libre bajo el término
de adaptación, que es la única fuerza que lo toca. La fuente siempre pasa por su
codificador congelado y el objetivo por el suyo.

\textbf{Lo que se mide acá es la fase 1 sola.} No hay dominio objetivo, no hay
adaptación, no hay segundo codificador: solo las 16 coordenadas que el codificador
de la fuente produce, y lo que un kernel gaussiano puede hacer con ellas.

\section{Qué muestra cada tabla}

\subsection{La geometría: direcciones efectivas y separación}

La covarianza de los rasgos de la fuente, $\Sigma_s = \tfrac{1}{n}\sum_i (z_i -
\bar z)(z_i - \bar z)^\top$, se descompone en
\[
\Sigma_s \;=\; \underbrace{\Sigma_{\text{dentro}}}_{\text{dispersión de cada clase}}
\;+\; \underbrace{\Sigma_{\text{entre}}}_{\text{separación entre clases}} .
\]
De ahí salen dos números. El primero es la \textbf{razón de participación} sobre
los autovalores $\lambda_1 \ge \dots \ge \lambda_{16}$ de $\Sigma_s$,
\[
  \mathrm{dim}_{\mathrm{ef}} \;=\; \frac{\left(\sum_k \lambda_k\right)^2}{\sum_k \lambda_k^2}
  \;\in\; [1, 16],
\]
que vale $1$ si toda la varianza está en una dirección y $16$ si es isótropa. Es
libre de escala: multiplicar los rasgos por cualquier constante no la mueve, así
que no confunde ``se encogió'' con ``se aplastó''. El segundo es
$\operatorname{tr}\Sigma_{\text{entre}} / \operatorname{tr}\Sigma_s$, la fracción
de la varianza que \emph{es} separación de clases.

\subsection{Lo que el kernel consume}
\noindent
\begin{minipage}[t]{0.49\textwidth}
  \centering
  \footnotesize\setlength{\tabcolsep}{3pt}
  \begin{tabular}{lrrrrr}
      \multicolumn{6}{c}{\textbf{SIN BatchNorm}} \\[2pt]
      \toprule
      familia & $\sigma$ & rango & sat. & $\bar K_{=}$ & $\bar K_{\neq}$ \\
      \midrule
      rollo & 49.2 & 45 & 25.7\% & 0.9667 & 0.5283 \\
      toroides & 23.6 & 97 & 26.0\% & 0.9674 & 0.4137 \\
      lunas & 19.4 & 1421 & 33.3\% & 0.9989 & 0.4190 \\
      \bottomrule
    \end{tabular}
\end{minipage}\hfill
\begin{minipage}[t]{0.49\textwidth}
  \centering
  \footnotesize\setlength{\tabcolsep}{3pt}
  \begin{tabular}{lrrrrr}
      \multicolumn{6}{c}{\textbf{CON BatchNorm}} \\[2pt]
      \toprule
      familia & $\sigma$ & rango & sat. & $\bar K_{=}$ & $\bar K_{\neq}$ \\
      \midrule
      rollo & 17.0 & 135 & 33.3\% & 0.9967 & 0.5722 \\
      toroides & 17.4 & 89 & 23.9\% & 0.9551 & 0.4794 \\
      lunas & 16.8 & 733 & 33.3\% & 0.9991 & 0.5763 \\
      \bottomrule
    \end{tabular}
\end{minipage}
\vspace{11pt}

\noindent
\begin{minipage}[t]{0.49\textwidth}
  \centering
  \footnotesize\setlength{\tabcolsep}{3pt}
  \begin{tabular}{lrrrrr}
      \multicolumn{6}{c}{\textbf{SIN BatchNorm, con KL}} \\[2pt]
      \toprule
      familia & $\sigma$ & rango & sat. & $\bar K_{=}$ & $\bar K_{\neq}$ \\
      \midrule
      rollo & 14.4 & 7 & 6.4\% & 0.9303 & 0.5723 \\
      toroides & 13.9 & 7 & 4.2\% & 0.9225 & 0.5603 \\
      lunas & 13.8 & 45 & 24.9\% & 0.9638 & 0.5692 \\
      \bottomrule
    \end{tabular}
\end{minipage}\hfill
\begin{minipage}[t]{0.49\textwidth}
  \centering
  \footnotesize\setlength{\tabcolsep}{3pt}
  \begin{tabular}{lrrrrr}
      \multicolumn{6}{c}{\textbf{CON BatchNorm, con KL}} \\[2pt]
      \toprule
      familia & $\sigma$ & rango & sat. & $\bar K_{=}$ & $\bar K_{\neq}$ \\
      \midrule
      rollo & 14.3 & 7 & 6.1\% & 0.9306 & 0.5712 \\
      toroides & 14.0 & 7 & 4.2\% & 0.9263 & 0.5657 \\
      lunas & 14.4 & 21 & 10.7\% & 0.9242 & 0.5637 \\
      \bottomrule
    \end{tabular}
\end{minipage}
\vspace{11pt}


El término usa un kernel gaussiano de instancia,
\[
  k(x,y) \;=\; \exp\!\left(-\frac{\lVert x - y\rVert^2}{2\sigma^2}\right),
\]
con \textbf{un solo $\sigma$} para los tres bloques $K^s$, $K^t$ y $K^{st}$ ---
porque los tres derivan del mismo kernel de instancia, y mezclar anchos rompe la
semidefinición positiva de la matriz conjunta, con lo cual $H_2$ dejaría de ser la
entropía de nada. Con \texttt{auto}, $\sigma$ es la mediana de las distancias por
pares. \textit{rango} es el cociente entre el percentil 99 y el 1 de esas
distancias: dice cuánto tendría que cubrir ese único $\sigma$. \textit{sat.} es el
porcentaje de pares con $k > 0{,}95$, que el kernel ya no distingue.

\textbf{Un hecho que conviene tener presente:} $k$ depende solo de $\lVert x-y
\rVert$, así que trasladar un dominio entero \emph{no mueve ninguna distancia
dentro de ese dominio}. Está medido: trasladar el objetivo por los $95{,}06$ de la
brecha cambia $K^t$ en $8{,}3 \times 10^{-6}$ y $K^{st}$ en $0{,}9995$. De los tres
bloques, uno solo ve un desplazamiento de medias.

\subsection{El contraste a nivel de clase, donde se ve la recta}
\noindent
\begin{minipage}[t]{0.49\textwidth}
  \centering
  \footnotesize\setlength{\tabcolsep}{3pt}
  \begin{tabular}{lrrr}
      \multicolumn{4}{c}{\textbf{SIN BatchNorm}} \\[2pt]
      \toprule
      familia & $K(0,1)$ & $K(0,2)$ & $K(1,2)$ \\
      \midrule
      rollo & 0.5863 & 0.5012 & 0.4974 \\
      toroides & 0.1386 & 0.5413 & 0.5613 \\
      lunas & 0.1222 & 0.6139 & 0.5208 \\
      \bottomrule
    \end{tabular}
\end{minipage}\hfill
\begin{minipage}[t]{0.49\textwidth}
  \centering
  \footnotesize\setlength{\tabcolsep}{3pt}
  \begin{tabular}{lrrr}
      \multicolumn{4}{c}{\textbf{CON BatchNorm}} \\[2pt]
      \toprule
      familia & $K(0,1)$ & $K(0,2)$ & $K(1,2)$ \\
      \midrule
      rollo & 0.5712 & 0.5799 & 0.5654 \\
      toroides & 0.4300 & 0.4732 & 0.5351 \\
      lunas & 0.5498 & 0.5860 & 0.5931 \\
      \bottomrule
    \end{tabular}
\end{minipage}
\vspace{11pt}

\noindent
\begin{minipage}[t]{0.49\textwidth}
  \centering
  \footnotesize\setlength{\tabcolsep}{3pt}
  \begin{tabular}{lrrr}
      \multicolumn{4}{c}{\textbf{SIN BatchNorm, con KL}} \\[2pt]
      \toprule
      familia & $K(0,1)$ & $K(0,2)$ & $K(1,2)$ \\
      \midrule
      rollo & 0.5676 & 0.5710 & 0.5784 \\
      toroides & 0.5587 & 0.5594 & 0.5628 \\
      lunas & 0.5694 & 0.5699 & 0.5682 \\
      \bottomrule
    \end{tabular}
\end{minipage}\hfill
\begin{minipage}[t]{0.49\textwidth}
  \centering
  \footnotesize\setlength{\tabcolsep}{3pt}
  \begin{tabular}{lrrr}
      \multicolumn{4}{c}{\textbf{CON BatchNorm, con KL}} \\[2pt]
      \toprule
      familia & $K(0,1)$ & $K(0,2)$ & $K(1,2)$ \\
      \midrule
      rollo & 0.5779 & 0.5698 & 0.5659 \\
      toroides & 0.5674 & 0.5664 & 0.5632 \\
      lunas & 0.5526 & 0.5675 & 0.5710 \\
      \bottomrule
    \end{tabular}
\end{minipage}
\vspace{11pt}


$K(a,b)$ es el valor medio del kernel entre la clase $a$ y la $b$. Esa matriz de
$3\times 3$ es, literalmente, \textbf{todo lo que un método de kernel sabe de las
clases}: los pesos por confianza de la Ec.~(24), la normalización por trazas de la
Ec.~(32), las cotas $L_c$ y $U_c$ de la Ec.~(35) y la mezcla marginal/condicional
leen todos esta misma matriz, ninguno la produce.

Con BatchNorm los tres pares quedan a menos de $0{,}09$ unos de otros: las tres
clases están aproximadamente equidistantes, un \textbf{triángulo en un plano}. Sin
BatchNorm, en toroides y lunas, un par queda en $0{,}12$--$0{,}14$ y los otros dos
en $0{,}52$--$0{,}61$, casi iguales entre sí. Esa es la firma de \textbf{tres
puntos ordenados sobre una recta}: los dos extremos lejos, y el del medio cerca de
ambos.

Y ahí está el problema, que es geométrico y no de ajuste. Si los tres centroides
son colineales, existe $t \in (0,1)$ con $\mu_2 = t\mu_0 + (1-t)\mu_1$, y entonces
la clase del medio \emph{no puede} estar lejos de las dos a la vez. El kernel la ve
parecida a ambas, y desde su punto de vista los extremos son indistinguibles. Para
que $C$ clases puedan ser mutuamente separables hacen falta al menos $C-1$
direcciones efectivas \textbf{antes} de calcular el kernel.

\section{¿Se arregla poniendo BatchNorm en una sola capa?}

La red es
\[
\text{Lin}(d,32) \;[\text{BN}]\; \text{ReLU} \;\;
\text{Lin}(32,32) \;[\text{BN}]\; \text{ReLU} \;\;
\text{Lin}(32,16) \;[\text{BN}] ,
\]
así que hay tres lugares posibles. El tercero, sobre la salida de 16 dimensiones,
es el único que normaliza \emph{exactamente lo que el kernel lee}, y no estaba en
la arquitectura: se agregó como brazo propio.

\begin{center}
\begin{tabular}{llrrrr}
  \toprule
   & & \multicolumn{2}{c}{\textbf{sin KL}} & \multicolumn{2}{c}{\textbf{con KL}} \\
  \cmidrule(lr){3-4}\cmidrule(lr){5-6}
  familia & BatchNorm en & dim.\,ef. & acierto & dim.\,ef. & acierto \\
  \midrule
    rollo & ninguna & 1.97 & 0.9401 $\pm$ 0.109 & 2.81 & 0.7355 $\pm$ 0.251 \\
     & 1.ª capa & 1.89 & 0.9355 $\pm$ 0.129 & 2.76 & 1.0000 $\pm$ 0.000 \\
     & 2.ª capa & 1.99 & 0.7321 $\pm$ 0.226 & 2.83 & 0.9846 $\pm$ 0.031 \\
     & salida & 2.01 & 0.4590 $\pm$ 0.165 & 3.11 & 0.9679 $\pm$ 0.064 \\
     & ambas ocultas & 1.97 & 0.6769 $\pm$ 0.014 & 2.80 & 0.9451 $\pm$ 0.079 \\
    \midrule
    toroides & ninguna & 1.11 & 0.3651 $\pm$ 0.030 & 2.83 & 0.2781 $\pm$ 0.068 \\
     & 1.ª capa & 1.68 & 0.7222 $\pm$ 0.055 & 2.79 & 0.7157 $\pm$ 0.139 \\
     & 2.ª capa & 1.90 & 0.4176 $\pm$ 0.088 & 2.85 & 0.3448 $\pm$ 0.223 \\
     & salida & 2.06 & 0.5830 $\pm$ 0.130 & 3.16 & 0.4046 $\pm$ 0.110 \\
     & ambas ocultas & 1.86 & 0.6299 $\pm$ 0.098 & 2.83 & 0.7657 $\pm$ 0.206 \\
    \midrule
    lunas & ninguna & 1.04 & 0.3306 $\pm$ 0.006 & 2.33 & 0.3346 $\pm$ 0.006 \\
     & 1.ª capa & 1.17 & 0.6534 $\pm$ 0.093 & 2.63 & 0.8398 $\pm$ 0.038 \\
     & 2.ª capa & 1.92 & 0.3731 $\pm$ 0.068 & 2.81 & 0.4080 $\pm$ 0.171 \\
     & salida & 2.00 & 0.2367 $\pm$ 0.057 & 2.80 & 0.2225 $\pm$ 0.106 \\
     & ambas ocultas & 1.95 & 0.4988 $\pm$ 0.133 & 2.83 & 0.7062 $\pm$ 0.109 \\
  \bottomrule
\end{tabular}
\end{center}

\textbf{Sí se arregla, y el lugar es la primera capa.} Sin KL, pasar de ninguna a
la primera lleva toroides de $0{,}3651$ a $\mathbf{0{,}7222}$ y lunas de $0{,}3306$
a $\mathbf{0{,}6534}$ --- por encima de ponerlo en \emph{las dos} capas ocultas,
que es la configuración con la que veníamos trabajando.

\textbf{Y la salida es el peor lugar de todos}, que es lo contrario de lo que
sugiere la intuición: normalizar justo los números que el kernel consume da
$0{,}4590$ en rollo y $0{,}2367$ en lunas, por debajo de no poner ninguno. Una
normalización sobre la salida fija la escala por coordenada \emph{después} de que
la representación ya se formó; no reparte la varianza entre direcciones, solo
reescala el aplastamiento que ya estaba. La primera capa, en cambio, actúa cuando
la representación todavía se está construyendo.

\section{Conclusión}

\begin{enumerate}\itemsep3pt
\item \textbf{Al kernel no le falta señal sin BatchNorm; le sobra.} El contraste
  $\bar K_{=} - \bar K_{\neq}$ es \emph{mayor} sin BatchNorm en toroides y lunas.
  La saturación es la misma en todos los brazos y casi nada muere. Así que $\sigma$,
  el rango dinámico y la saturación --- los sospechosos de siempre --- no son la
  diferencia.

\item \textbf{La diferencia es la dimensionalidad efectiva.} Las dos familias que
  fracasan sin BatchNorm ponen el $94{,}8\%$ y el $97{,}9\%$ de la varianza en una
  sola dirección, con $1{,}11$ y $1{,}04$ direcciones efectivas. La que funciona
  conserva $1{,}97$, las mismas que cualquier brazo con BatchNorm.

\item \textbf{La palanca no es BatchNorm, es la dimensionalidad.} BatchNorm es
  \emph{una} forma de conseguirla; rollo sin él conserva $58/41$ y llega a
  $0{,}9401$. El término KL es otra: sus anclas son ortonormales, o sea le
  prescribe $C$ direcciones por construcción, y sube las direcciones efectivas a
  $\approx 2{,}8$ en todas las configuraciones medidas.

\item \textbf{Ningún arreglo río abajo alcanza.} $\sigma$, los pesos por confianza,
  la mezcla marginal/condicional y la normalización por cotas leen todos la misma
  matriz de $3\times3$. Si esa matriz dice que la clase del medio está entre las
  otras dos, ninguno de ellos puede decir otra cosa.

\item \textbf{La mejor configuración medida, por familia, 5 semillas:}
  \begin{center}
  \begin{tabular}{llr}
    \toprule
    familia & configuración & acierto \\
    \midrule
        rollo & BatchNorm en la primera capa, con KL & 1.0000 $\pm$ 0.000 \\
    toroides & BatchNorm en ambas capas ocultas, con KL & 0.7657 $\pm$ 0.206 \\
    lunas & BatchNorm en la primera capa, con KL & 0.8398 $\pm$ 0.038 \\
    \bottomrule
  \end{tabular}
  \end{center}
  contra $0{,}6769$, $0{,}6299$ y $0{,}4988$ del punto de partida (BatchNorm en
  ambas capas ocultas, sin KL).
\end{enumerate}

\vspace{6pt}
\begin{center}\small\itshape
Todo número de este documento sale de una corrida medida sobre 5 semillas; el
azar con 3 clases es $0{,}3333$. Las familias son de juguete y nada dice todavía
que estos valores transfieran a los datos reales.
\end{center}

\end{document}
